\section{同样平均光子数，为什么电子相干性仍不同}
\label{chap:feqo-photon-statistics}

把电子初态取为 \(|0\rangle_e\)，光的初态取任意单模密度算符 \(\rho_\gamma\)。在短时间或弱耦合下设 \(\xi=\int g(t)\dd t\)，并要求 \(|\xi|\sqrt{\langle N_\gamma\rangle+1}\ll1\)。式~\eqref{eq:feqo-quantum-pinem-hamiltonian} 的一阶传播为
\[
U=I-\ii\xi B^\dagger a-\ii\xi^*B a^\dagger+O(|\xi|^2).
\]
为求到 \(|\xi|^2\) 阶的人口，还需把二阶返回原态的振幅计入；利用幺正性确定其总和即可。记 \(\bar n=\operatorname{Tr}(\rho_\gamma N_\gamma)\)，则
\begin{equation}
P_{+1}=|\xi|^2\bar n+O(|\xi|^4),\quad
P_{-1}=|\xi|^2(\bar n+1)+O(|\xi|^4),\quad
P_0=1-|\xi|^2(2\bar n+1)+O(|\xi|^4).
\label{eq:feqo-weak-quantum-populations}
\end{equation}
这里 \(a^\dagger a=N_\gamma\) 给吸收权重，\(aa^\dagger=N_\gamma+1\) 给发射权重；额外的 1 是真空涨落。若有初末态失谐，\(\xi\) 应换成相应的有限时间 Fourier 积分，不能只积 \(g\)。

为什么更高边带能够区分统计？在相位固定且忽略反冲的相互作用中，第二次吸收来自传播指数的平方项 \(-\xi^2(B^\dagger a)^2/2\)。初始光子数为 \(n\) 时，\(a^2|n\rangle=\sqrt{n(n-1)}|n-2\rangle\)，于是对任意光初态，领先的双吸收概率为
\begin{equation}
P_{+2}=\frac{|\xi|^4}{4}
\langle N_\gamma(N_\gamma-1)\rangle+O(|\xi|^6).
\label{eq:feqo-double-absorption-statistics}
\end{equation}
这个式子的适用还要求相关高阶光子数矩有限，且 \(|\xi|\sqrt n\) 在初态有显著概率的 \(n\) 范围内小。固定平均数 \(\bar n=1\) 时，单光子数态给 \(\langle N(N-1)\rangle=0\)，相干态给 1，热态给 2；因此三者的一阶边带可以相同，而第二吸收边带依次为 \(0,|\xi|^4/4,|\xi|^4/2\) 的领先阶。若电子开始就占有多个边带、反冲使两次吸收失谐不同，须重新计算有序时间积分，不能继续把 \(\xi^2/2\) 当成普适振幅。

相同的 \(\bar n\) 在这个最低非零阶给相同人口，因而只比较弱场谱峰不能区分光统计。电子相干项却已经不同。按式~\eqref{eq:feqo-partial-trace-components} 对光求迹，
\begin{equation}
(\rho_e)_{1,0}=-\ii\xi\langle a\rangle+O(|\xi|^2),\qquad
(\rho_e)_{-1,0}=-\ii\xi^*\langle a^\dagger\rangle+O(|\xi|^2).
\label{eq:feqo-quantum-coherence-first-order}
\end{equation}
Fock 态 \(|n\rangle\) 与热态都满足 \(\langle a\rangle=0\)，所以一阶能量相干为零；相干态 \(|\alpha\rangle\) 满足 \(\langle a\rangle=\alpha\)，于是保留相位。比较时可选 \(n=1\)、\(|\alpha|^2=1\)、热平均数 \(\bar n=1\)，并固定同一个 \(\xi\)。人口先相同，电子相干先不同；更高阶人口还会读出 \(\langle N_\gamma(N_\gamma-1)\rangle\) 等统计矩。

对纯初始电子，保留式~\eqref{eq:feqo-weak-quantum-populations} 与式~\eqref{eq:feqo-quantum-coherence-first-order} 的最低非零阶，纯度为
\[
\operatorname{Tr}\rho_e^2
=1-2|\xi|^2\bigl(2\bar n+1-2|\langle a\rangle|^2\bigr)+O(|\xi|^3).
\]
Fock 和热态的括号是 \(2\bar n+1\)，相干态则是 1。该式只在扰动区内比较纯度；高阶项和具体脉冲相位可能改变更细的差别。

相干态为何会恢复上一章的经典场，可用位移算符看清。令 \(D(\alpha)=e^{\alpha a^\dagger-\alpha^*a}\)，则 \(D^\dagger aD=a+\alpha\)。在位移后的表象中，式~\eqref{eq:feqo-quantum-pinem-hamiltonian} 变为确定驱动 \(\hbar g\alpha B^\dagger+\mathrm{h.c.}\) 加上量子涨落耦合 \(\hbar gB^\dagger a+\mathrm{h.c.}\)。令 \(|\alpha|\to\infty\)、\(g\to0\)，同时固定 \(g\alpha\) 和作用时间，确定驱动不变而相对涨落尺度 \(1/|\alpha|\to0\)。此时电子演化趋向经典场的相位调制；仅说“光子数大”而不说明固定什么量，不构成一个受控极限。

式~\eqref{eq:feqo-joint-displacement} 还允许全阶比较三种光输入，而不在 \(|\xi|\sqrt{\bar n+1}\) 处截断。固定交换总量 \(N_\gamma+L_e=n\)，以末态光子数 \(m\geq0\) 标记联合基 \(|n-m,m\rangle\)。在这条链上，\(A\) 把 \(m\) 降一，\(A^\dagger\) 把 \(m\) 升一，作用系数与普通谐振子完全相同。定义位移矩阵元 \(D_{mn}(\zeta):=\langle m|e^{\zeta a^\dagger-\zeta^*a}|n\rangle\)。由式~\eqref{eq:feqo-joint-displacement} 的两个指数分别展开，先降 \(n-k\) 次、再升 \(m-k\) 次，得到有限和
\begin{equation}
D_{mn}(\zeta)=e^{-\mu/2}\sqrt{m!n!}
\sum_{k=0}^{\min(m,n)}
\frac{\zeta^{m-k}(-\zeta^*)^{n-k}}
{k!(m-k)!(n-k)!},\qquad \mu=|\zeta|^2.
\label{eq:feqo-displacement-matrix}
\end{equation}
这一式不要求读者事先认识关联 Laguerre 多项式。对入射 \(|0,n\rangle\)，末态振幅恰为 \(D_{mn}|n-m,m\rangle\)，故电子第 \(\ell\) 条边带的概率是 \(P_\ell^{(n)}=|D_{n-\ell,n}|^2\)（\(n-\ell\geq0\)）；其它 \(\ell\) 的概率为零。不同 \(\ell\) 留下不同光子数，所以电子约化态在能量基对角。

单光子 \(n=1\) 的前几项可以完全写开：
\[
P_{+1}^{(1)}=\mu e^{-\mu},\qquad
P_0^{(1)}=(1-\mu)^2e^{-\mu},\qquad
P_{-1}^{(1)}=\frac{\mu(2-\mu)^2}{2}e^{-\mu}.
\]
\(\mu=1\) 时中心边带恰被相消抽空，但概率由其它所有边带承接；只保留这三项不能强制它们的和等于一。若初态是真空 \(n=0\)，式~\eqref{eq:feqo-displacement-matrix} 退化为式~\eqref{eq:feqo-vacuum-joint-state} 的泊松分布。这个可直接检验的极限也固定了所有阶乘和符号。

相干光初态 \(|\alpha\rangle\) 的数态系数为 \(c_n=e^{-\bar n/2}\alpha^n/\sqrt{n!}\)，其中 \(\bar n=|\alpha|^2\)。同一末态光子数 \(m\) 可由不同初始 \(n=m+\ell\) 到达，所以联合振幅和电子约化矩阵分别是
\begin{equation}
C_{\ell m}=c_{m+\ell}D_{m,m+\ell}(\zeta),\qquad
(\rho_e)_{\ell\ell'}=\sum_{m\geq\max(0,-\ell,-\ell')}
C_{\ell m}C_{\ell'm}^* .
\label{eq:feqo-coherent-reduced-exact}
\end{equation}
负指标的 \(c_n\) 定义为零。这个和式保留来自不同初始光子数的相位干涉。热光初态则是对角混合，权重 \(p_n^{\rm th}=\bar n^n/(1+\bar n)^{n+1}\)，因此
\(P_\ell^{\rm th}=\sum_{n\geq\max(0,\ell)}p_n^{\rm th}|D_{n-\ell,n}|^2\)，电子非对角元为零。固定 \(\bar n=1\) 后，Fock 态 \(n=1\)、相干态 \(|\alpha|=1\)、热态可在完全相同的 \(\zeta\) 下比较，而不是把平均输入能量的差别误认为光统计。

这些无穷和也有透明的数值误差。若只保留初始光子数 \(n\leq N\)，任一电子人口的遗漏不超过输入光分布的尾概率 \(\sum_{n>N}p_n\)，因为每个固定 \(n\) 的条件概率不超过一。\(\bar n=1\) 的热态尾和正好是 \(2^{-(N+1)}\)；相干态尾和是 \(1-e^{-1}\sum_{n=0}^N1/n!\)。在计算前选择 \(N\) 使这个数小于目标误差，便有真正受控的截断。

\begin{exercise}
由式~\eqref{eq:feqo-displacement-matrix} 求 \(D_{0,1},D_{1,1},D_{2,1}\)，再用单光子结果解释 \(P_0^{(1)}\) 在 \(\mu=1\) 消失的原因。
\end{exercise}
